\section{Código fuente}\label{apx:codigo}
\justifying

Este apéndice presenta \emph{fragmentos curados} del código fuente que ilustran las decisiones de diseño más importantes. El código fuente completo (backend, frontend, tests, scripts y fixtures) se entrega en el archivo \texttt{codigo\_fuente.zip} de la carpeta de entrega, generado con \texttt{git archive HEAD}.

\subsection{Estructura del repositorio}

\begin{Verbatim}[fontsize=\small, frame=single]
proyecto_terminal_cyad/
|-- backend/
|   |-- config/                # settings, urls, wsgi
|   |-- core/                  # utilidades: pagination, exceptions
|   |-- accounts/              # Usuario, Profesor, JWT
|   |-- catalogos/             # Departamento, Licenciatura, UEA, Periodo
|   |-- documentos/            # CartaTematica, RequisitoRecuperacion
|   |-- autoevaluacion/        # Formulario, Pregunta, Respuesta, Nivel
|   |-- reportes/              # Dashboards y cumplimiento
|   |-- scripts/               # seed_users, seed_demo, create_db
|   |-- tests/                 # ~4200 lineas de tests con pytest
|   |-- manage.py
|   `-- requirements.txt
|-- frontend/
|   |-- src/
|   |   |-- api/               # clientes HTTP (auth, documentos, ...)
|   |   |-- auth/              # AuthContext, ProtectedRoute, RoleRoute
|   |   |-- components/ui/     # Button, Alert, Badge, Table, Modal
|   |   |-- layouts/           # AdminLayout, ProfesorLayout
|   |   |-- pages/
|   |   |   |-- admin/         # dashboards y CRUD por catalogo
|   |   |   |-- profesor/      # cartas, requisitos, autoevaluacion
|   |   |   `-- publico/       # vista Explorar y documentos publicados
|   |   |-- App.jsx            # rutas + guards
|   |   `-- main.jsx           # bootstrap React + TanStack Query
|   |-- index.html
|   |-- package.json
|   `-- vite.config.js
|-- docs/
|   |-- reporte/               # este documento
|   `-- manuales/              # instalacion y usuario
`-- README.md
\end{Verbatim}

\subsection{Backend --- Modelo de Usuario con roles}

Fragmento de \texttt{backend/accounts/models.py}: modelo de usuario personalizado con dos roles.

\begin{Verbatim}[fontsize=\small, frame=single]
class Usuario(AbstractBaseUser, PermissionsMixin):
    class Rol(models.TextChoices):
        ADMIN    = "ADMIN",    "Administrador"
        PROFESOR = "PROFESOR", "Profesor"

    email     = models.EmailField(unique=True)
    nombre    = models.CharField(max_length=150)
    rol       = models.CharField(max_length=10, choices=Rol.choices)
    is_active = models.BooleanField(default=True)
    is_staff  = models.BooleanField(default=False)

    USERNAME_FIELD = "email"
    REQUIRED_FIELDS = ["nombre"]

    objects = UsuarioManager()
\end{Verbatim}

\subsection{Backend --- JWT y throttling}

Fragmento de \texttt{backend/config/settings.py}.

\begin{Verbatim}[fontsize=\small, frame=single]
REST_FRAMEWORK = {
    "DEFAULT_AUTHENTICATION_CLASSES": (
        "rest_framework_simplejwt.authentication.JWTAuthentication",
    ),
    "DEFAULT_PERMISSION_CLASSES": (
        "rest_framework.permissions.IsAuthenticated",
    ),
    "DEFAULT_THROTTLE_CLASSES": (
        "rest_framework.throttling.ScopedRateThrottle",
    ),
    "DEFAULT_THROTTLE_RATES": {
        "login": "5/min",
    },
    "PAGE_SIZE": 20,
}

SIMPLE_JWT = {
    "ACCESS_TOKEN_LIFETIME":  timedelta(minutes=60),
    "REFRESH_TOKEN_LIFETIME": timedelta(days=7),
    "ROTATE_REFRESH_TOKENS":  True,
    "AUTH_HEADER_TYPES":      ("Bearer",),
}
\end{Verbatim}

\subsection{Backend --- Documento academico base}

Fragmento de \texttt{backend/documentos/models.py}. Clase abstracta con campos comunes y \emph{snapshot} de identidad del profesor.

\begin{Verbatim}[fontsize=\small, frame=single]
class DocumentoAcademicoBase(TimeStampedModel):
    class Estado(models.TextChoices):
        BORRADOR   = "BORRADOR",   "Borrador"
        PUBLICADO  = "PUBLICADO",  "Publicado"

    profesor         = models.ForeignKey(
        Profesor, on_delete=models.SET_NULL, null=True, blank=True
    )
    profesor_nombre  = models.CharField(max_length=200, blank=True)
    profesor_correo  = models.EmailField(blank=True)

    uea       = models.ForeignKey(UEA,     on_delete=models.PROTECT)
    periodo   = models.ForeignKey(Periodo, on_delete=models.PROTECT)
    nombre_grupo = models.CharField(max_length=50)
    id_grupo     = models.CharField(max_length=50)
    horario      = models.CharField(max_length=200, blank=True)
    modalidad    = models.CharField(
        max_length=15,
        choices=[("PRESENCIAL","Presencial"),("REMOTO","Remoto"),("MIXTO","Mixto")],
    )
    estado    = models.CharField(
        max_length=15,
        choices=Estado.choices,
        default=Estado.BORRADOR,
    )

    class Meta:
        abstract = True
\end{Verbatim}

\subsection{Backend --- Scoring ponderado de autoevaluacion}\label{lst:scoring}

Fragmento simplificado de \texttt{backend/autoevaluacion/views.py}: calculo del puntaje ponderado por seccion, con manejo por tipo de pregunta y \emph{snapshot} del peso en \texttt{RespuestaSeccion}.

\begin{Verbatim}[fontsize=\small, frame=single]
def _calcular_puntaje(respuesta):
    resultado_secciones = []
    item_by_pregunta = {r.pregunta_id: r for r in respuesta.items.all()}
    for seccion in respuesta.formulario.secciones.all():
        sec_obt, sec_max = Decimal("0"), Decimal("0")
        for pregunta in seccion.preguntas.all():
            if pregunta.tipo in Pregunta.TIPOS_NO_PUNTABLES:
                continue
            item = item_by_pregunta.get(pregunta.id)
            if pregunta.tipo == Pregunta.Tipo.CUADRICULA:
                opts, filas = list(pregunta.opciones.all()), list(pregunta.filas.all())
                if opts and filas:
                    sec_max += max(o.puntos for o in opts) * len(filas)
                if item:
                    sec_obt += sum(c.opcion.puntos for c in item.celdas.all())
            elif pregunta.tipo == Pregunta.Tipo.ESCALA_LINEAL:
                cfg = pregunta.config or {}
                factor = Decimal(str(cfg.get("puntos_factor", 1)))
                sec_max += Decimal(str(cfg.get("max", 5))) * factor
                if item and item.valor_texto:
                    sec_obt += Decimal(item.valor_texto) * factor
            # ... SI_NO, CASILLAS, OPCION_UNICA / LISTA_DESPLEGABLE ...
        porcentaje = (sec_obt / sec_max * 100).quantize(Decimal("0.01")) \
                     if sec_max > 0 else Decimal("0")
        resultado_secciones.append({
            "seccion": seccion, "obtenido": sec_obt, "maximo": sec_max,
            "porcentaje": porcentaje, "peso": seccion.peso,
        })

    porcentaje_global = sum(
        s["porcentaje"] * s["peso"] / Decimal("100") for s in resultado_secciones
    ).quantize(Decimal("0.01"))

    with transaction.atomic():
        respuesta.porcentaje = porcentaje_global
        respuesta.nivel_desempeno = NivelDesempeno.objects.filter(
            formulario=respuesta.formulario,
            porcentaje_min__lte=porcentaje_global,
            porcentaje_max__gte=porcentaje_global,
        ).first()
        respuesta.save()
        for s in resultado_secciones:
            RespuestaSeccion.objects.update_or_create(
                respuesta=respuesta, seccion=s["seccion"],
                defaults={
                    "puntaje":        s["obtenido"],
                    "puntaje_maximo": s["maximo"],
                    "porcentaje":     s["porcentaje"],
                    "peso":           s["peso"],   # snapshot para desglose historico
                },
            )
\end{Verbatim}

\subsection{Backend --- Permisos por rol}\label{lst:permissions}

Fragmento del sistema de permisos que garantiza que cada profesor solo edite sus documentos.

\begin{Verbatim}[fontsize=\small, frame=single]
class IsOwnerProfesorOrAdminReadOnly(permissions.BasePermission):
    def has_permission(self, request, view):
        return request.user.is_authenticated

    def has_object_permission(self, request, view, obj):
        if request.user.rol == Usuario.Rol.ADMIN:
            return request.method in permissions.SAFE_METHODS
        return obj.profesor and obj.profesor.usuario_id == request.user.id
\end{Verbatim}

\subsection{Frontend --- Cliente HTTP con refresh automatico de JWT}\label{lst:jwt-refresh}

Fragmento de \texttt{frontend/src/api/client.js}.

\begin{Verbatim}[fontsize=\small, frame=single]
const client = axios.create({
  baseURL: import.meta.env.VITE_API_URL,
});

client.interceptors.request.use((config) => {
  const token = sessionStore.getAccessToken();
  if (token) config.headers.Authorization = `Bearer ${token}`;
  return config;
});

let refreshing = null;

client.interceptors.response.use(
  (r) => r,
  async (error) => {
    const { config, response } = error;
    if (response?.status === 401 && !config.__retried) {
      config.__retried = true;
      refreshing = refreshing ?? refreshAccessToken();
      const newToken = await refreshing.finally(() => (refreshing = null));
      if (newToken) {
        config.headers.Authorization = `Bearer ${newToken}`;
        return client(config);
      }
    }
    return Promise.reject(error);
  }
);
\end{Verbatim}

\subsection{Frontend --- Guard de rutas por rol}\label{lst:role-route}

Fragmento de \texttt{frontend/src/auth/RoleRoute.jsx}.

\begin{Verbatim}[fontsize=\small, frame=single]
export function RoleRoute({ role, children }) {
  const { user, isLoading } = useAuth();
  if (isLoading) return <Loading />;
  if (!user) return <Navigate to="/login" replace />;
  if (user.rol !== role) return <Navigate to="/" replace />;
  return children;
}

export const AdminRoute    = ({ children }) =>
  <RoleRoute role="ADMIN"    >{children}</RoleRoute>;
export const ProfesorRoute = ({ children }) =>
  <RoleRoute role="PROFESOR" >{children}</RoleRoute>;
\end{Verbatim}

\subsection{Frontend --- Renderizador dinamico de preguntas}\label{lst:renderer_preguntas}

Fragmento de \code{frontend/src/pages/profesor/autoevaluacion/AutoevaluacionFormPage.jsx} que despacha por tipo de pregunta.

\begin{Verbatim}[fontsize=\small, frame=single]
function PreguntaRenderer({ pregunta, value, onChange }) {
  switch (pregunta.tipo) {
    case "TEXTO_CORTO":       return <TextInput  {...{value, onChange}} />;
    case "TEXTO_LARGO":       return <TextArea   {...{value, onChange}} />;
    case "OPCION_UNICA":      return <RadioGroup opciones={pregunta.opciones} {...{value, onChange}} />;
    case "CASILLAS":          return <CheckGroup opciones={pregunta.opciones} {...{value, onChange}} />;
    case "SI_NO":             return <YesNo      {...{value, onChange}} />;
    case "ESCALA_LINEAL":     return <LinearScale config={pregunta.config} {...{value, onChange}} />;
    case "LISTA_DESPLEGABLE": return <Select     opciones={pregunta.opciones} {...{value, onChange}} />;
    case "CUADRICULA":        return <Grid       config={pregunta.config} filas={pregunta.filas} {...{value, onChange}} />;
    default:                  return <UnknownType tipo={pregunta.tipo} />;
  }
}
\end{Verbatim}

\subsection{Backend --- Comando de importacion CSV para UEA}\label{lst:import_csv}

Comando de gestion \texttt{python manage.py cargar\_catalogos\_csv} para carga masiva idempotente. Fragmento simplificado.

\begin{Verbatim}[fontsize=\small, frame=single]
class Command(BaseCommand):
    help = "Carga catalogos (Areas, Licenciaturas, Posgrados, UEAs) desde CSV"

    def add_arguments(self, parser):
        parser.add_argument("--ueas", type=str, help="Path al CSV de UEAs")

    def handle(self, *args, **opts):
        ruta = opts["ueas"]
        with open(ruta, newline="", encoding="utf-8") as fh:
            for fila in csv.DictReader(fh):
                UEA.objects.update_or_create(
                    clave=fila["clave"],
                    licenciatura=Licenciatura.objects.get(clave=fila["licenciatura"]),
                    defaults={
                        "nombre":    fila["nombre"],
                        "trimestre": int(fila["trimestre"]),
                        "creditos":  int(fila["creditos"]),
                        "tipo":      fila["tipo"],
                    },
                )
\end{Verbatim}

\subsection{Backend --- Constraint XOR y unicidad parcial de UEA}\label{lst:uea_constraints}

Fragmento de \texttt{backend/catalogos/models.py}: reglas de integridad declarativas que fuerzan que una UEA pertenezca a licenciatura \emph{o} a posgrado (nunca a ambos ni a ninguno), y que la clave sea unica dentro de cada programa.

\begin{Verbatim}[fontsize=\small, frame=single]
class Meta:
    constraints = [
        CheckConstraint(
            condition=(
                (Q(licenciatura__isnull=False) & Q(posgrado__isnull=True))
                | (Q(licenciatura__isnull=True)  & Q(posgrado__isnull=False))
            ),
            name="uea_xor_licenciatura_posgrado",
        ),
        UniqueConstraint(
            fields=["clave", "licenciatura"],
            condition=Q(licenciatura__isnull=False),
            name="uea_clave_unica_por_licenciatura",
        ),
        UniqueConstraint(
            fields=["clave", "posgrado"],
            condition=Q(posgrado__isnull=False),
            name="uea_clave_unica_por_posgrado",
        ),
    ]
\end{Verbatim}

\subsection{Backend --- Ciclo de vida de Formulario}\label{lst:formulario_states}

Fragmento de \texttt{backend/autoevaluacion/models.py}: validacion de estructura y transiciones de la maquina de estados \texttt{BORRADOR} $\leftrightarrow$ \texttt{PUBLICADO} $\leftrightarrow$ \texttt{CERRADO}.

\begin{Verbatim}[fontsize=\small, frame=single]
def _validar_estructura_para_publicar(self):
    secciones = list(self.secciones.all())
    if not secciones:
        raise ValueError("Debe definir al menos una seccion antes de publicar.")
    suma = sum(s.peso for s in secciones)
    if abs(suma - Decimal("100")) > Decimal("0.01"):
        raise ValueError(f"La suma de pesos debe ser 100% (actual: {suma}%).")
    if self.preguntas.filter(seccion__isnull=True).exists():
        raise ValueError("Hay preguntas sin seccion asignada.")

def publicar(self):
    self._validar_estructura_para_publicar()
    if self.estado == self.Estado.CERRADO and self.respuestas.exists():
        self.version += 1     # publicar_revision incrementa version
    self.estado = self.Estado.PUBLICADO
    self.published_at = timezone.now()
    self.save()

def cerrar(self):
    self.estado = self.Estado.CERRADO
    self.closed_at = timezone.now()
    self.save()

def reabrir(self):
    self.estado = self.Estado.PUBLICADO
    self.save()

def despublicar(self):
    # Vuelve a BORRADOR sin cambiar version ni borrar respuestas.
    self.estado = self.Estado.BORRADOR
    self.save()

def publicar_revision(self):
    # CERRADO -> PUBLICADO incrementando version.
    self._validar_estructura_para_publicar()
    self.version += 1
    self.estado = self.Estado.PUBLICADO
    self.published_at = timezone.now()
    self.save()
\end{Verbatim}

\subsection{Backend --- Login con distincion de cuenta desactivada}\label{lst:login_diff}

Fragmento de \texttt{backend/accounts/serializers.py}. El \emph{serializer} personalizado distingue \emph{credenciales invalidas} de \emph{cuenta desactivada} devolviendo un codigo \texttt{account\_disabled} sin filtrar la existencia del correo.

\begin{Verbatim}[fontsize=\small, frame=single]
class CustomTokenObtainPairSerializer(TokenObtainPairSerializer):
    def validate(self, attrs):
        email = (attrs.get(self.username_field) or "").strip().lower()
        password = attrs.get("password") or ""
        usuario = Usuario.objects.filter(email__iexact=email).first()
        if usuario and usuario.check_password(password) and not usuario.is_active:
            raise AuthenticationFailed(
                {"code": "account_disabled",
                 "detail": "Tu cuenta se encuentra desactivada. "
                           "Contacta al administrador para reactivarla."},
                code="account_disabled",
            )
        data = super().validate(attrs)   # credenciales OK, usuario activo
        u = self.user
        data["user"] = {"id": u.id, "email": u.email,
                        "nombre": u.nombre, "rol": u.rol}
        return data
\end{Verbatim}

\subsection{Frontend --- Generacion de XLSX por periodo}\label{lst:xlsx_por_periodo}

Fragmento simplificado de \code{frontend/src/pages/admin/reportes/ReportesPage.jsx}: descarga en el navegador con \emph{SheetJS}, agrupando las filas por periodo en hojas separadas del mismo libro.

\begin{Verbatim}[fontsize=\small, frame=single]
async function descargarDocumento(publicPath, fetcher) {
  const params = { estado: "PUBLICADO", page_size: 500 }
  if (periodoId) params.periodo = periodoId
  if (deptoId)   params.profesor__departamento = deptoId

  const res  = await fetcher(params)
  const rows = res.data?.results ?? res.data ?? []

  const porPeriodo = new Map()
  for (const row of rows) {
    const clave = row.periodo_clave ?? "Sin periodo"
    if (!porPeriodo.has(clave)) porPeriodo.set(clave, [])
    porPeriodo.get(clave).push(row)
  }

  const wb = XLSX.utils.book_new()
  if (porPeriodo.size === 0) {
    XLSX.utils.book_append_sheet(wb, XLSX.utils.json_to_sheet([]), "Sin datos")
  } else {
    for (const [clave, groupRows] of porPeriodo) {
      const ws = XLSX.utils.json_to_sheet(
        groupRows.map((r) => rowToExcel(r, publicPath))
      )
      XLSX.utils.book_append_sheet(wb, ws, sheetName(clave))
    }
  }
  XLSX.writeFile(wb, `${publicPath}-${Date.now()}.xlsx`)
}
\end{Verbatim}

\subsection{Nota sobre el listado completo}

Los archivos aquí incluidos son una selección representativa de los fragmentos más ilustrativos. El código fuente completo (backend, frontend, \emph{tests}, scripts y fixtures) se entrega en el archivo \texttt{codigo\_fuente.zip}. Para regenerarlo desde el repositorio se ejecuta desde la raíz del proyecto:

\begin{quote}
\ttfamily git archive HEAD -o docs/reporte/entregables/codigo\_fuente.zip
\end{quote}

Adicionalmente, la documentación viva de la API se genera automáticamente vía \emph{drf-spectacular} y se sirve en \texttt{/api/docs/} (interfaz Swagger UI) y \texttt{/api/schema/} (esquema OpenAPI en YAML).
